% =====================================================================
%  INTEGRATION PLAN -- three proposals into two IEEE S&P papers
%
%  Standalone. Set as Overleaf's main document; pdfLaTeX + BibTeX.
%    pdflatex INTEGRATION_PLAN && bibtex INTEGRATION_PLAN \
%      && pdflatex INTEGRATION_PLAN && pdflatex INTEGRATION_PLAN
% =====================================================================
\documentclass[11pt,a4paper]{article}
% =====================================================================
%  Shared preamble for the three research proposals.
%
%  Each proposal \input{}s this file and is otherwise standalone, so any
%  one of them can be set as Overleaf's main document and compiled on its
%  own. Nothing here is specific to a single proposal.
% =====================================================================
\usepackage[T1]{fontenc}
\usepackage[utf8]{inputenc}
% Times, matching the parent manuscript. mathptmx is in every TeX Live and on
% Overleaf; lmodern is deliberately not used, since it is absent from minimal
% distributions and this preamble should compile anywhere.
\usepackage{mathptmx}
\usepackage[a4paper,margin=25mm]{geometry}
\usepackage{amsmath,amssymb,amsthm}
\usepackage{booktabs}
\usepackage{tabularx}
\usepackage{xcolor}
\usepackage{titlesec}
\usepackage{enumitem}
\usepackage{microtype}
\usepackage[hidelinks,breaklinks]{hyperref}

% ---- shared style tokens, matching the parent manuscript ------------
\definecolor{netcol}{RGB}{31,78,121}
\definecolor{autocol}{RGB}{18,120,90}
\definecolor{bridgecol}{RGB}{176,96,0}
\definecolor{accent}{RGB}{140,20,20}

\newcommand{\layernet}{\textcolor{netcol}{\textbf{network}}}
\newcommand{\layerauto}{\textcolor{autocol}{\textbf{autonomy}}}
\newcommand{\dataset}[1]{\textsc{#1}}
\newcommand{\code}[1]{\texttt{#1}}
\newcommand{\Real}{\mathbb{R}}
\newcommand{\Norm}[1]{\left\lVert #1\right\rVert}
\newcommand{\MCR}{\mathrm{MCR}}

\newtheorem{claim}{Claim}
\newtheorem{proposition}{Proposition}
\newtheorem{definition}{Definition}
\theoremstyle{remark}
\newtheorem{remark}{Remark}

% ---- section styling -------------------------------------------------
\titleformat{\section}{\large\bfseries\color{netcol}}{\thesection}{0.7em}{}
\titleformat{\subsection}{\normalsize\bfseries}{\thesubsection}{0.6em}{}
\titlespacing*{\section}{0pt}{1.4em}{0.6em}

\setlist[itemize]{leftmargin=1.35em,topsep=0.3em,itemsep=0.25em}
\setlist[enumerate]{leftmargin=1.6em,topsep=0.3em,itemsep=0.25em}

% ---- a boxed "what this rests on" callout ---------------------------
\newenvironment{honestly}
  {\par\medskip\noindent\begin{tabular}{@{}p{2.5mm}@{\hspace{3mm}}p{0.9\linewidth}@{}}
   \textcolor{accent}{\rule{2.5mm}{\baselineskip}} &
   \small\itshape}
  {\end{tabular}\par\medskip}

% ---- the proposal header --------------------------------------------
\newcommand{\proposalheader}[3]{%
  \begin{center}
  {\Large\bfseries #1\par}\medskip
  {\large\itshape #2\par}\medskip
  {\small #3\par}
  \end{center}
  \medskip\hrule\medskip}

\usepackage{array}
\newcolumntype{Y}{>{\raggedright\arraybackslash}X}
\newcolumntype{P}[1]{>{\raggedright\arraybackslash}p{#1}}
\newcommand{\core}{\textcolor{autocol}{\textbf{I-core}}}
\newcommand{\appx}{\textcolor{netcol}{\textbf{I-appx}}}
\newcommand{\sok}{\textcolor{bridgecol}{\textbf{II}}}
\newcommand{\drop}{\textbf{drop}}

\begin{document}

\proposalheader
  {Integration Plan}
  {Three Follow-on Proposals, Two IEEE S\&P Papers}
  {RobustUAVs.ai research programme $\cdot$ Roger Nick Anaedevha $\cdot$
   2026-10-01\\
   Supersedes the standalone framing of P1, P2 and P3; those files are kept
   unchanged as the record of the original proposals.}

% =====================================================================
\section{The decision}
% =====================================================================
\textbf{Two papers, not one and not three.}

\begin{itemize}
\item \textbf{Paper I} (S\&P main track) \textbf{merges P1 and P2}.
  \emph{What Authenticity Costs in Metres: A Certified Budget for Cryptography
  and Sensing in UAV Navigation.}
\item \textbf{Paper II} (S\&P \emph{Systematization of Knowledge} track)
  \textbf{is P3, broadened}.
  \emph{SoK: The Quantifier Gap Between Certified Robustness and Airworthiness
  Assurance.}
\end{itemize}

\noindent
The reason is a single observation that only becomes visible when P1 and P2 are
read side by side. The parent manuscript's certificate is a budget,
\begin{equation}
\label{eq:budget}
\rho \;=\; \Delta \cdot \gamma \cdot e^{L T_c} \;\le\; m ,
\end{equation}
and P1 and P2 each attack one of its factors. P1 says $\gamma$ depends on which
navigation source is trusted, and that \emph{authenticating} a source is what
makes it trustworthy. P2 says $\Delta$ does not care why an input is stale, and
that \emph{authenticating} a message makes it stale. The same operation
therefore appears on both sides of Eq.~\eqref{eq:budget} with opposite signs:
authentication is simultaneously a credit (lower $\gamma$) and a debit (higher
$\Delta$). Neither proposal can say which wins. Together they can, and the
answer is the paper.

% =====================================================================
\section{Why not three papers}
% =====================================================================
Each of P1 and P2 is weaker alone than its share of the merge, and each says so
in its own risk section.

\begin{itemize}
\item \textbf{P1 alone is a measurement paper.} Its central claim is described
  in its own text as ``arithmetic rather than new theory''; its contribution is
  the per-modality $\gamma$. That is a solid journal contribution
  (\emph{TIFS}, \emph{TCPS}, which is where P1 targeted it) and a weak S\&P one.
\item \textbf{P2 alone risks being read as post-quantum benchmarking.} Its risk
  section names exactly this: ``reviewers may read it as post-quantum
  benchmarking''. Its mitigation, ``lead with the budget result'', is only
  available if there is something for the budget to trade \emph{against}. P1
  supplies it.
\item \textbf{Merged, there is a falsifiable result neither has}: a crossover
  inequality deciding when authentication increases the certified mission floor
  and when it decreases it (Paper~I, Proposition~1). That is a finding about
  security engineering, not a measurement of a component.
\end{itemize}

% =====================================================================
\section{Why not one paper}
% =====================================================================
P3 does not merge, and forcing it in would damage both halves.

\begin{itemize}
\item \textbf{Different contribution type.} P1 and P2 produce measurements, a
  theorem corollary and an attack. P3 produces a mapping, a gap analysis and a
  generator. S\&P main-track reviewers judge the first kind; the second kind is
  what the SoK track exists for.
\item \textbf{Different audience.} P3's own framing: it ``changes who the
  result is \emph{for}'', from a security audience to a certification audience.
  A paper cannot change its audience halfway through.
\item \textbf{No shared measurement.} P3 reuses the parent's committed results
  and needs nothing from P1 or P2. Bundling it would make Paper~II's schedule
  hostage to Paper~I's hardware campaign for no gain.
\item \textbf{Page budget.} Paper~I already carries two measurement campaigns in
  13 pages. P3's mapping tables alone would take three.
\end{itemize}

\begin{honestly}
P3 as written is not yet an S\&P SoK. It maps \emph{one} certificate (ours) onto
standards. An SoK must systematise a \emph{literature}. Paper~II therefore
broadens P3 from ``our certificate as evidence'' to ``every class of certified
guarantee for learned autonomy, organised by what it quantifies over'' --- see
Paper~II, \S2. Without that broadening it should go to DASC, which P3 already
named as a venue, rather than to S\&P.
\end{honestly}

% =====================================================================
\section{Element-by-element categorisation}
% =====================================================================
Every element of the three proposals, and where it goes. Legend:
\core{} = Paper~I main body;
\appx{} = Paper~I appendix;
\sok{} = Paper~II;
\drop{} = not carried forward.

\subsection{From P1 (certified PNT substitution)}
\begin{center}\small
\begin{tabularx}{\linewidth}{@{}P{0.36\linewidth}P{0.09\linewidth}Y@{}}
\toprule
\textbf{Element} & \textbf{Goes to} & \textbf{Why} \\
\midrule
Claim: $\gamma$ is a property of the navigation source, not the attack &
\core & The ``credit'' side of the budget. Without it P2 has nothing to trade
against. \\
Claim: authentication discharges the re-anchoring hypothesis & \core &
Splits Paper~I's analysis into two regimes (delay vs.\ persistent denial); in
the denial regime it is what makes \emph{any} certificate non-vacuous. \\
Claim: the certificate induces a sensor-selection rule & \core{} (rule)
\appx{} (optimality proof) & The rule is short and is the operator-facing
output; its brute-force verification is evidence, its proof is length. \\
$\gamma$ for GNSS & \core & Done. \textbf{Use $1.625$\,m/s} (campaign supremum),
not P1's $1.195$/$1.365$, which predate the interface campaign. \\
$\gamma$ inertial-only & \core & Data already on disk; about a week. Do first. \\
$\gamma$ visual--inertial & \core & Public VIO datasets; second measured
modality beyond GNSS. \\
$\gamma$ LEO signals of opportunity & \appx & No receiver logs. A
\code{physics\_model} value from published RMSE is honest only if tagged and
kept out of the headline. \\
Fusion model $\gamma_{\mathrm{eff}}=\min_s\gamma_s$ & \core & Conservative and
unglamorous, as P1 asked. Filter-weighted variant to the appendix. \\
Denial-duration sweep & \core & Provides the regime-(ii) evidence. \\
Schema change (\code{pnt\_source}, \code{pnt\_authenticated}) & \core{} (one
paragraph) & Additive, keeps every existing record valid. \\
PNT Source Selector page & platform & Merged into one Paper~I page; see \S7. \\
\bottomrule
\end{tabularx}
\end{center}

\subsection{From P2 (the security tax)}
\begin{center}\small
\begin{tabularx}{\linewidth}{@{}P{0.36\linewidth}P{0.09\linewidth}Y@{}}
\toprule
\textbf{Element} & \textbf{Goes to} & \textbf{Why} \\
\midrule
Budget $\rho=(\Delta+\Delta_{\mathrm{auth}})\gamma e^{LT_c}\le m$ & \core &
Becomes the merged budget, with $\gamma\to\gamma_{\mathrm{eff}}(S,a)$. \\
Prop.: detector tightness buys cryptographic headroom & \core &
Kept, and sharpened: Paper~I shows it holds for \emph{affordability} while the
opposite holds for \emph{benefit} (Paper~I, Proposition~2). \\
Prop.: a margin below which nothing certifies & \core{} (already partly shown) &
At the parent's constants, $m=2$\,m certifies nothing \emph{even with zero
authentication cost} (the budget is $-0.12$\,s). The proposition is partly a
consequence of existing results; the new part is naming the margin per
transport. \\
Prop.: ranking by size is the wrong ranking & \core & The ``cliff'' result:
fragmentation and bus occupancy, not per-message latency, bind. \\
Measure sign/verify latency on MCU and companion board & \core & The
measurement that keeps the result from restating its inputs. \\
DroneCAN fragmentation model & \core & Validated against benign HCRL timing. \\
Induced-verification attack & \core & The S\&P hook. Must be positioned
against TESLA-era prior art on signature-flooding DoS~\cite{perrig2000tesla}:
the mechanism is known, the mission-level bound is new. \\
TESLA-style delayed authentication as a scheme & \core{} (new) & Not in P2.
It is a scheme whose $\Delta_{\mathrm{auth}}$ is an explicit, tunable
parameter --- the cleanest possible test of the budget. \\
Candidate scheme set pinned to FIPS 203/204/205 & \core & As P2 said: lets a
future scheme be placed without re-running. \\
Crypto Budget page & platform & Merged into one Paper~I page; see \S7. \\
\code{pq\_swarm\_kyber} (MDPI) & \drop & Key agreement, not per-message
authentication; not load-bearing for Paper~I, and the supervisor's guidance is
to avoid MDPI unless core. \\
\bottomrule
\end{tabularx}
\end{center}

\subsection{From P3 (certificate as compliance artifact)}
\begin{center}\small
\begin{tabularx}{\linewidth}{@{}P{0.36\linewidth}P{0.09\linewidth}Y@{}}
\toprule
\textbf{Element} & \textbf{Goes to} & \textbf{Why} \\
\midrule
``Certified radii are per-input; assurance objectives are per-operation'' &
\sok{} (central axis) & This sentence, buried in P3's honesty box, is the
organising idea of the SoK. Promoted to the title. \\
Vocabulary mapping (radius, tube, operating region $\leftrightarrow$ ODD,
contingency volume) & \sok & One row per certificate class, not just ours. \\
Discharge / evidence / untouched gap analysis & \sok & Applied to every
surveyed guarantee class, not only the four in the parent paper. \\
Simplex runtime assurance as the bridge & \sok & Developed as P3 asked. \\
SORA Annex A evidence generator & \sok{} (artifact) & What makes the SoK more
than a survey. \\
Compliance Pack export & platform & One export action. \\
\bottomrule
\end{tabularx}
\end{center}

% =====================================================================
\section{Categorisation by contribution type}
% =====================================================================
The same elements, grouped by the kind of claim a reviewer will judge them as.

\begin{center}\small
\begin{tabularx}{\linewidth}{@{}P{0.17\linewidth}YY@{}}
\toprule
\textbf{Type} & \textbf{Paper I} & \textbf{Paper II} \\
\midrule
Theory & Merged budget; crossover inequality; two-regime split
(delay vs.\ denial); selection rule &
Quantifier taxonomy (per-input, per-trajectory, per-mission, per-population) \\
Measurement & $\gamma$ for GNSS, inertial, VIO; $\Delta_{\mathrm{auth}}$ for
HMAC, Ed25519, ML-DSA, SLH-DSA, TESLA; bus occupancy & none (reuses
committed results) \\
Attack & Induced verification, with bound and mitigation & none \\
Systematisation & --- & Literature of certified guarantees for learned
autonomy, mapped onto SORA / ED-324 / EASA learning assurance \\
Artifact & Budget calculator, crossover map, selection rule & SORA Annex A
evidence generator \\
\bottomrule
\end{tabularx}
\end{center}

% =====================================================================
\section{Which of the parent's limitations each paper closes}
% =====================================================================
\begin{center}\small
\begin{tabularx}{\linewidth}{@{}P{0.36\linewidth}P{0.13\linewidth}P{0.13\linewidth}Y@{}}
\toprule
\textbf{Parent limitation} & \textbf{Paper I} & \textbf{Paper II} & \textbf{Note} \\
\midrule
Re-anchoring is assumed & \textbf{closes} & --- & Authentication makes the
hypothesis satisfiable; denial regime analysed explicitly. \\
$\gamma$ calibrated on one platform, one regime & \textbf{partly} & --- &
Adds modalities; still \emph{not} a cruise campaign. That gap remains
data-gated (\code{PENDING\_ON\_DATA.md}). \\
One attack surface (mesh delay) & \textbf{partly} & --- & Adds the bus
(\code{intra\_bus}) and C2 link (\code{c2\_link}) on the defence side. \\
Guarantee bounds a population, not a flight & --- & \textbf{reframes} & The
SoK's central finding is that this is a general property of the field, not a
weakness of one paper. \\
Evidence-class imbalance & no & no & Unchanged by either. Say so. \\
\bottomrule
\end{tabularx}
\end{center}

% =====================================================================
\section{Platform integration}
% =====================================================================
Three proposed pages collapse to two.
\begin{itemize}
\item \textbf{Certified Budget} (Paper~I): replaces both the \emph{PNT Source
  Selector} and the \emph{Crypto Budget} pages. Choose a sensor suite with
  authentication states, a transport, a scheme, a message rate and a margin;
  see the budget, the crossover sign, and the bus occupancy, live. One
  endpoint, \code{/api/certify/budget}.
\item \textbf{Compliance Pack} (Paper~II): an export action, unchanged from P3.
\end{itemize}

% =====================================================================
\section{Readiness against the S\&P bar}
% =====================================================================
\begin{center}\small
\begin{tabularx}{\linewidth}{@{}P{0.28\linewidth}YY@{}}
\toprule
\textbf{Reviewer will ask} & \textbf{Paper I} & \textbf{Paper II} \\
\midrule
Is there a new result, not just a measurement? & Yes: the crossover and the
cliff. & Yes, if broadened: the quantifier gap as a field-wide finding. \\
Is the threat model explicit? & Inherited from parent; adds the
induced-verification adversary. & Not applicable; must state its inclusion
criteria instead. \\
Are numbers measured or modelled? & Mixed; every number carries a provenance
tag. LEO is modelled. & Reuses measured results. \\
Is there an attack? & Yes, with a bound and a mitigation. & No, and an SoK
does not need one. \\
Biggest risk & Hardware access for the MCU benchmarks. & A thin literature
corpus. Needs a systematic search with stated criteria (\S2 of Paper~II). \\
Ready to start & Now: inertial $\gamma$ and the budget arithmetic need nothing
new. & Now: needs no measurement. \\
\bottomrule
\end{tabularx}
\end{center}

% =====================================================================
\section{Sequencing}
% =====================================================================
\begin{enumerate}
\item \textbf{Week 1--2, Paper I:} inertial-only $\gamma$ from the Whelan logs
  already on disk. This is also the measurement that most strengthens the
  \emph{parent} manuscript, whose cruise discussion currently asserts this
  value from an error budget.
\item \textbf{Week 1--6, Paper II:} systematic literature search; this runs
  in parallel and depends on nothing.
\item \textbf{Week 3--10, Paper I:} MCU and companion-board benchmarks; VIO
  $\gamma$; fragmentation validation against HCRL.
\item \textbf{Week 8--12, Paper I:} induced-verification attack and mitigation.
\item \textbf{Submit Paper II first.} It is cheaper and lower-risk, and an SoK
  that cites the parent establishes the vocabulary Paper~I then uses.
\end{enumerate}

\begin{honestly}
Paper I carries the larger risk, because it has two measurement campaigns and
either can fail. If the MCU benchmarks cannot be run on real flight-controller
hardware, the fallback is the pinned-ARM measurement P2 already named, tagged as
such --- and the crossover result survives, because it depends on the
\emph{ratio} $\Delta_{\mathrm{auth}}/\Delta(\theta)$, which a pinned ARM target
still bounds. If the $\gamma$ campaign yields only GNSS and inertial, two
modalities still establish that $\gamma$ is source-dependent. The paper degrades
gracefully; it does not collapse to P1 or P2 alone.
\end{honestly}

\small
\bibliographystyle{plain}
\bibliography{proposals}

\end{document}
